\documentclass[12pt,a4paper]{article}

\usepackage{amsmath,amssymb}
\usepackage{graphicx}
\usepackage{booktabs}
\usepackage{multirow}
\usepackage{siunitx}
\usepackage{hyperref}
\usepackage[numbers]{natbib}
\usepackage{lineno}
\usepackage{xcolor}
\usepackage[margin=2.5cm]{geometry}
\usepackage{authblk}
\usepackage{longtable}
\usepackage{array}

\modulolinenumbers[5]

\begin{document}

\title{Improved Lattice Parameter Prediction of Cubic High-Entropy Alloys\\
via Structure-Specific DFT Volume Size Factors\\
and Element-Wise Additive Decomposition}

\author[1]{Satoshi Minamoto}
\affil[1]{National Institute for Materials Science (NIMS), Tsukuba, Ibaraki, Japan}

\date{}
\maketitle

\begin{abstract}
\noindent
Accurate prediction of lattice parameters is essential for computational
screening of high-entropy alloys (HEAs). Alonso's model, which applies
King's (1966) experimental volume size factors ($\Omega_\text{sf}$) to
correct Vegard's law, achieves an RMSE of 0.023~\AA{} on 64 cubic HEAs.
We improve upon this by (i)~deriving \emph{structure-specific}
$\Omega_\text{sf}$ from first-principles calculations of 3,511 binary
intermetallic compounds, using B2 data for BCC HEAs and L1$_2$ data for
FCC HEAs; (ii)~introducing element-wise additive decomposition
$\Omega_\text{sf}(A{-}B) \approx \delta_A + \delta_B$, which replaces
220 pairwise parameters with 39 element-level parameters while preserving
prediction accuracy; and (iii)~quantifying for the first time why the
conventional atomic size mismatch $\delta r$ cannot capture structural
information, whereas volume-derived quantities inherently do.
Our best model achieves RMSE = 0.0210~\AA{} (8.6\% improvement over
Alonso), with FCC HEAs predicted at 0.0096~\AA{}---approaching the
experimental noise floor ($\sigma \approx 0.016$~\AA).  Validation on 20
independent HEAs confirms generalisation (RMSE = 0.020~\AA).  The
element-wise $\delta$ parameters provide a compact, periodic-table-style
descriptor that predicts $\Omega_\text{sf}$ for any element pair without
additional DFT calculations, enabling rapid HEA screening.

\medskip
\noindent\textbf{Keywords:}
High-entropy alloys, Lattice parameter prediction,
Volume size factor, First-principles calculation,
Additive decomposition, Atomic size mismatch
\end{abstract}

\linenumbers

%% =====================================================================
\section{Introduction}
\label{sec:intro}
%% =====================================================================

High-entropy alloys (HEAs) represent a paradigm shift in materials design,
consisting of five or more principal elements in near-equimolar ratios that
form simple solid-solution phases~\cite{Yeh2004,Cantor2004}. These alloys
exhibit exceptional combinations of mechanical properties, including high
strength, ductility, fracture toughness, and corrosion
resistance~\cite{Miracle2017}. The vast compositional space---spanning
millions of possible element combinations---makes computational screening
essential for accelerating discovery~\cite{Senkov2015}.

The lattice parameter $a$ is a fundamental structural property that
determines atomic packing, elastic moduli, and thermal expansion.
Vegard's law~\cite{Vegard1921} provides the simplest estimate:
\begin{equation}
  V_\text{Vegard} = \sum_i c_i V_i, \qquad
  a = \left(n_\text{auc}\,V_\text{Vegard}\right)^{1/3},
  \label{eq:vegard}
\end{equation}
where $c_i$ is the atomic fraction of element $i$, $V_i$ the elemental
atomic volume, and $n_\text{auc}$ the number of atoms per unit cell
(2 for BCC, 4 for FCC).

King~\cite{King1966} introduced the volume size factor $\Omega_\text{sf}$
to quantify the fractional volume deviation when a solute atom dissolves
into a host lattice. Alonso~\cite{Alonso2021} applied this concept to
multicomponent HEAs:
\begin{equation}
  V = n_\text{auc} \left[
    \sum_i c_i V_i + \sum_i c_i V_i \sum_{j \ne i} c_j\,
    \Omega_\text{sf}(j, i)
  \right],
  \label{eq:alonso_eq10}
\end{equation}
where $\Omega_\text{sf}(j,i)$ is King's volume size factor for the $j$--$i$
pair. This pairwise correction captures chemical bonding effects beyond
Vegard's law and achieves RMSE $\approx 0.023$~\AA{} on 64 cubic HEAs.

Despite its success, Alonso's approach has several limitations:
(1)~the $\Omega_\text{sf}$ values originate from King's 1966 experimental
data, which may be unavailable or inaccurate for modern HEA-relevant pairs;
(2)~the same $\Omega_\text{sf}$ is used regardless of crystal structure,
despite the distinct atomic environments in BCC versus FCC;
(3)~for novel element combinations, missing pairwise data forces fallback
to pure Vegard's law.

The conventional atomic size mismatch parameter
\begin{equation}
  \delta r = 100 \sqrt{\sum_i c_i\!\left(1 - \frac{r_i}{\bar{r}}\right)^{\!2}},
  \qquad \bar{r} = \sum_i c_i\,r_i,
  \label{eq:delta_r}
\end{equation}
where $r_i$ are pure-element metallic radii, is widely used for
single-phase prediction~\cite{Zhang2008,Guo2011}. However, as we
demonstrate rigorously in this work, $\delta r$ is a function of
$(c_i, r_i)$ alone and contains no crystallographic structure information
whatsoever.

In this work, we address these limitations through four innovations:
\begin{enumerate}
  \item We derive $\Omega_\text{sf}$ from \emph{first-principles} DFT
    calculations of 3,511 binary compounds from the Materials Project
    and OQMD, with structure-specific separation (B2 for BCC, L1$_2$
    for FCC).
  \item We introduce \emph{element-wise additive decomposition}
    $\Omega_\text{sf}(A{-}B) \approx \delta_A + \delta_B$, reducing
    220 pairwise parameters to 39 element-level parameters.
  \item We prove mathematically and validate computationally (905
    binary pairs) that $\delta r$ cannot absorb structural information,
    whereas volume-derived radii from DFT inherently capture
    structure-dependent effects.
  \item We validate on 20 independent HEAs not present in the Alonso
    training set, confirming generalisation.
\end{enumerate}


%% =====================================================================
\section{Methodology}
\label{sec:method}
%% =====================================================================

\subsection{Data Sources}
\label{sec:data}

\subsubsection{Binary Compound Data}
Binary intermetallic compounds were collected from three sources:
\begin{itemize}
  \item \textbf{Materials Project (MP)}~\cite{Jain2013}: B2 and L1$_2$
    compounds computed with PAW-PBE.
  \item \textbf{OQMD}~\cite{Saal2013}: Additional B2 and L1$_2$ entries.
  \item \textbf{VASP calculations (this work)}: 1,483 B2 and 2,808 L1$_2$
    compounds computed with PAW-PBE (ENCUT = 520~eV,
    $k$-mesh $\geq 12\times12\times12$, EDIFF = $10^{-6}$~eV).
\end{itemize}
In total, 9,182 binary compound entries were obtained. After filtering
for convergence, physical lattice constants (2.0--8.0~\AA), and
requiring at least two entries per element pair, the dataset yields
pairwise $\Omega_\text{sf}$ for 156 B2 and 101 L1$_2$ element pairs
(220 unique pairs total). For the comprehensive structural analysis
(Section~\ref{sec:structural}), 905 element pairs have data for all
three structures (L1$_2$ $A_3B$, $B_3A$, and B2 $AB$).

\subsubsection{HEA Validation Data}
\label{sec:hea_data}
The primary validation set consists of 64 cubic HEAs (29 BCC, 35 FCC)
from Alonso~\cite{Alonso2021} Table~2 ($a = 2.883$--$3.856$~\AA, 22
elements). An independent test set of 20 HEAs (8 BCC, 12 FCC) was
collected from separate
publications~\cite{Senkov2013,Yao2016,Otto2013,Niu2017,Wang2019,
Zaddach2013,Zhang2017nat} not included in Alonso's dataset.

\subsubsection{Multiphase HEA Database}
\label{sec:multiphase_data}
For phase stability analysis, 54 HEAs were compiled from seven
literature sources~\cite{Zhang2008,Guo2011,Yang2012,Senkov2013mp},
classified as solid solution (SS, $N = 25$), intermetallic-containing
(IM, $N = 23$), or amorphous (AM, $N = 6$).

\subsection{Volume Size Factor Computation}
\label{sec:omega}

For each binary compound $A_x B_y$ with DFT lattice constant
$a_\text{DFT}$, the fractional volume deviation from Vegard's law is:
\begin{equation}
  \Omega_\text{sf}(A{-}B) =
  \frac{V_\text{DFT} - V_\text{Vegard}}{V_\text{Vegard}},
  \label{eq:omega_def}
\end{equation}
where $V_\text{DFT} = a_\text{DFT}^3 / Z$ ($Z = 2$ for B2, $Z = 4$
for L1$_2$) and $V_\text{Vegard} = \sum_i x_i V_i$ using King's
atomic volumes~\cite{King1966}. The median across all compounds for
each pair mitigates outlier effects.

Structure-specific factors are maintained:
\begin{align}
  \Omega_\text{sf}^\text{B2}(A{-}B) &= \text{median}\left\{
    \Omega_\text{sf} \mid \text{B2 compounds of } A{-}B
  \right\}, \label{eq:omega_b2} \\
  \Omega_\text{sf}^{\text{L1}_2}(A{-}B) &= \text{median}\left\{
    \Omega_\text{sf} \mid \text{L1}_2 \text{ compounds of } A{-}B
  \right\}. \label{eq:omega_l12}
\end{align}

\subsection{Structure-Specific Lattice Prediction}
\label{sec:ss_eq10}

We generalise Eq.~\eqref{eq:alonso_eq10} with a structure-dependent
scaling factor $\gamma_s$:
\begin{equation}
  V = n_\text{auc} \left[
    \sum_i c_i V_i + \gamma_s \sum_i \sum_{j \ne i}
    c_i\, c_j\, V_j\, \Omega_\text{sf}^{(s)}(i, j)
  \right],
  \label{eq:scaled_eq10}
\end{equation}
where $s \in \{\text{BCC}, \text{FCC}\}$, with
$\Omega_\text{sf}^{(\text{BCC})} = \Omega_\text{sf}^\text{B2}$ and
$\Omega_\text{sf}^{(\text{FCC})} = \Omega_\text{sf}^{\text{L1}_2}$.
Grid search over $\gamma \in [-0.5, 2.0]$ yields optimal values
$\gamma_\text{BCC} = 1.45$ and $\gamma_\text{FCC} = 1.08$.

\subsection{Element-Wise Additive Decomposition}
\label{sec:additive}

To overcome the missing-pair problem and provide a compact descriptor,
we decompose the pairwise $\Omega_\text{sf}$ into element-level
contributions:
\begin{equation}
  \Omega_\text{sf}(A{-}B) \approx \delta_A + \delta_B,
  \label{eq:additive}
\end{equation}
where $\delta_A$ is a scalar parameter assigned to element $A$.
This is solved by ordinary least squares on the overdetermined system
($P$ pairs, $N$ elements, $P \gg N$):
\begin{equation}
  \min_{\boldsymbol{\delta}} \sum_{(A,B)}
    \left[\Omega_\text{sf}(A{-}B) - \delta_A - \delta_B\right]^2.
  \label{eq:additive_lsq}
\end{equation}
The decomposition is performed separately for B2 and L1$_2$ data,
yielding $\delta_A^{\text{B2}}$ and $\delta_A^{\text{L1}_2}$.

The additive model replaces the pairwise $\Omega_\text{sf}$
correction in Eq.~\eqref{eq:scaled_eq10}:
\begin{equation}
  V \approx n_\text{auc} \left[
    \sum_i c_i V_i + \gamma_s \sum_i \sum_{j \ne i}
    c_i\, c_j\, V_j\, \left(\delta_i^{(s)} + \delta_j^{(s)}\right)
  \right].
  \label{eq:additive_pred}
\end{equation}
This enables prediction for \emph{any} element pair where both
elements have $\delta$ values, regardless of whether their pairwise
$\Omega_\text{sf}$ has been computed.

\subsection{Structure-Dependent Effective Radii}
\label{sec:eff_radii}

To investigate whether effective radii absorb structural information,
we derive radii from DFT atomic volumes. For L1$_2$ $A_3B$, the
per-atom volume is modelled as:
\begin{equation}
  V_\text{per atom}(A_3B) = 0.75\,V_\text{maj}(A) +
    0.25\,V_\text{min}(B),
  \label{eq:l12_volume}
\end{equation}
where $V_\text{maj}$ and $V_\text{min}$ are effective atomic volumes
for the majority (face-centred) and minority (corner) sites,
respectively. A joint least-squares optimisation with $2N_\text{el}$
unknowns is performed simultaneously over all L1$_2$ data. The
effective radii are then:
\begin{equation}
  r_\text{eff} = \left(\frac{3V_\text{eff}}{4\pi}\right)^{1/3}.
  \label{eq:r_eff}
\end{equation}

\subsection{Packing-Constrained Radius Analysis}
\label{sec:packing}

We also evaluate whether DFT-derived packing-constrained radii can
capture structural information. The nearest-neighbour distance is:
\begin{align}
  d_\text{nn}^{\text{L1}_2} &= a / \sqrt{2}, \label{eq:dnn_l12} \\
  d_\text{nn}^{\text{B2}}   &= a\sqrt{3} / 2. \label{eq:dnn_b2}
\end{align}
Under the touching-sphere constraint, $r_A + r_B = d_\text{nn}$.
A single set of radii $(r_A, r_B)$ is fitted to all structures
simultaneously by least squares.

\subsection{Machine Learning Residual Correction}
\label{sec:ml}

To explore whether nonlinear corrections improve upon the physics
model, we train ML models on the residuals
$\Delta a_i = a_i^\text{exp} - a_i^\text{Eq.\ref{eq:scaled_eq10}}$.
Seven algorithms are evaluated via leave-one-out cross-validation
(LOO-CV): Ridge, XGBoost, GPR, SVR, Random Forest, Cubist, and
ensemble combinations. A 23-dimensional feature vector is constructed
including Vegard lattice constant, DFT Eq.~10 prediction, composition
statistics, $\delta r$, VEC, and $\Omega_\text{sf}$ pair statistics.

\subsection{Noise Floor Estimation}
\label{sec:noise}

Three pairs of duplicate HEA compositions in the validation set yield
an estimated measurement noise of $\sigma \approx 0.016$~\AA{} (from
$E[\text{range}] = 1.128\,\sigma$ for $n = 2$ Gaussian samples),
establishing a fundamental lower bound on achievable RMSE.


%% =====================================================================
\section{Results}
\label{sec:results}
%% =====================================================================

\subsection{HEA Lattice Parameter Prediction}
\label{sec:results_hea}

Table~\ref{tab:comparison} summarises the prediction accuracy on the
64 validation HEAs.

\begin{table}[htbp]
\centering
\caption{Lattice parameter prediction accuracy for 64 cubic HEAs.}
\label{tab:comparison}
\begin{tabular}{@{}lcccc@{}}
\toprule
Method & RMSE (\AA) & MAE (\AA) & MAPE (\%) & $R^2$ \\
\midrule
\multicolumn{5}{l}{\textit{Baseline methods}} \\
Alonso Vegard       & 0.0280 & 0.0167 & 0.504 & 0.9845 \\
Alonso Eq.~10 ($\gamma = 1$) & 0.0229 & 0.0141 & 0.418 & 0.9897 \\
\midrule
\multicolumn{5}{l}{\textit{This work: pairwise $\Omega_\text{sf}$}} \\
King Vegard         & 0.0227 & 0.0139 & 0.421 & 0.9899 \\
DFT Eq.~10 ($\gamma = 1$) & 0.0268 & 0.0202 & 0.588 & 0.9859 \\
DFT Eq.~10 SS ($\gamma_\text{BCC}, \gamma_\text{FCC}$)
  & \textbf{0.0210} & \textbf{0.0117} & \textbf{0.348} & \textbf{0.9914} \\
\midrule
\multicolumn{5}{l}{\textit{This work: additive decomposition}} \\
Additive $\delta$ + $\gamma$ (Approach~1) & 0.0215 & --- & --- & --- \\
\midrule
\multicolumn{5}{l}{\textit{ML residual corrections (LOO-CV)}} \\
SS Eq.~10 + Ridge   & 0.0212 & 0.0116 & 0.344 & 0.9912 \\
SS Eq.~10 + XGBoost & 0.0212 & 0.0116 & 0.345 & 0.9911 \\
SS Eq.~10 + GPR     & 0.0214 & 0.0124 & 0.366 & 0.9910 \\
SS Eq.~10 + RF      & 0.0221 & 0.0117 & 0.349 & 0.9904 \\
4-way Ensemble      & 0.0210 & 0.0115 & 0.343 & 0.9914 \\
\midrule
Noise floor ($\sigma$) & \multicolumn{4}{c}{$\approx 0.016$~\AA} \\
\bottomrule
\end{tabular}
\end{table}

The structure-specific DFT Eq.~10 (``DFT Eq.~10 SS'') achieves the best
physics-model RMSE of 0.0210~\AA, an 8.6\% improvement over Alonso's
Eq.~10 (0.0229~\AA). ML residual corrections (Ridge, XGBoost, GPR)
provide only marginal improvement ($< 0.002$~\AA), confirming that
the physics model captures the dominant variability.

\subsection{BCC/FCC Breakdown}
\label{sec:bcc_fcc}

\begin{table}[htbp]
\centering
\caption{BCC/FCC breakdown of RMSE (\AA) for key methods.}
\label{tab:bcc_fcc}
\begin{tabular}{@{}lccc@{}}
\toprule
Method & BCC ($N\!=\!29$) & FCC ($N\!=\!35$) & Overall \\
\midrule
Alonso Eq.~10           & 0.0309 & 0.0132 & 0.0229 \\
King Vegard             & 0.0328 & 0.0085 & 0.0227 \\
DFT Eq.~10 SS           & 0.0293 & 0.0096 & 0.0210 \\
Additive $\delta$ + $\gamma$ & 0.0310 & 0.0084 & 0.0215 \\
\bottomrule
\end{tabular}
\end{table}

The FCC subset is predicted with exceptional accuracy (RMSE = 0.0096~\AA),
below the noise floor. The BCC predictions (0.0293~\AA) improve upon
Alonso by 5.2\%. The disparity reflects: (a)~FCC HEAs being dominated
by 3$d$ transition metals with similar volumes; (b)~BCC HEAs spanning
wider atomic size ranges (Al through Zr/Hf); (c)~B2 compounds being
ordered intermetallics that differ structurally from random BCC solid
solutions.

\subsection{Independent Test Validation}
\label{sec:indep}

Table~\ref{tab:indep} shows performance on 20 HEAs not present in the
Alonso training set.

\begin{table}[htbp]
\centering
\caption{Independent test set validation (20 HEAs).}
\label{tab:indep}
\begin{tabular}{@{}lccc@{}}
\toprule
Method & RMSE (\AA) & BCC (8) & FCC (12) \\
\midrule
Vegard           & 0.0202 & 0.0260 & 0.0152 \\
Alonso Eq.~10    & 0.0261 & 0.0345 & 0.0186 \\
DFT Eq.~10 SS    & \textbf{0.0200} & \textbf{0.0271} & \textbf{0.0134} \\
\bottomrule
\end{tabular}
\end{table}

The independent-test RMSE of 0.0200~\AA{} is comparable to the training
RMSE (0.0210~\AA), confirming generalisation. The FCC independent-test
RMSE (0.0134~\AA) is close to the noise floor.

\subsection{Element-Wise Additive Decomposition Parameters}
\label{sec:delta_table}

The pairwise $\Omega_\text{sf}$ is decomposed into element-level
parameters $\delta_A$ via Eq.~\eqref{eq:additive_lsq}.
Table~\ref{tab:delta} presents the full parameter set.

\begin{longtable}{@{}l
  S[table-format=+1.4] S[table-format=+1.4]
  S[table-format=+1.4] S[table-format=+1.4] S[table-format=+1.4]@{}}
\caption{Element-wise additive decomposition parameters $\delta_A$ for
  B2 and L1$_2$ structures, and structure-dependent effective radii
  from DFT joint optimisation (51 elements). $r_\text{pure}$ is the
  King metallic radius; $r_\text{maj}$ and $r_\text{min}$ are L1$_2$
  majority and minority effective radii; $\Delta r = r_\text{eff} -
  r_\text{pure}$.}
\label{tab:delta} \\
\toprule
{Element} & {$\delta^\text{B2}$} & {$\delta^{\text{L1}_2}$} &
{$r_\text{pure}$ (\AA)} & {$\Delta r_\text{maj}$ (\AA)} &
{$\Delta r_\text{min}$ (\AA)} \\
\midrule
\endfirsthead
\toprule
{Element} & {$\delta^\text{B2}$} & {$\delta^{\text{L1}_2}$} &
{$r_\text{pure}$ (\AA)} & {$\Delta r_\text{maj}$ (\AA)} &
{$\Delta r_\text{min}$ (\AA)} \\
\midrule
\endhead
\bottomrule
\endfoot
Ag &  -0.0086 &  +0.0237 & 1.597 & +0.016 & +0.058 \\
Al &  -0.0489 &  -0.0187 & 1.583 & -0.015 & +0.060 \\
As &   {---}  &   {---}  & 1.722 & -0.074 & +0.062 \\
Au &  -0.0021 &  +0.0191 & 1.594 & +0.029 & +0.019 \\
B  &   {---}  &   {---}  & 1.200 & -0.562 & +0.167 \\
Ba &   {---}  &   {---}  & 2.285 & -0.071 & -0.104 \\
Be &  -0.0144 &   {---}  & 1.246 & -0.015 & +0.112 \\
Bi &   {---}  &   {---}  & 2.037 & -0.068 & -0.074 \\
Ca &  -0.1578 &  -0.1064 & 2.184 & -0.152 & -0.215 \\
Cd &   {---}  &   {---}  & 1.727 & -0.030 & -0.047 \\
Ce &  +0.0492 &  -0.0447 & 1.857 & +0.043 & +0.119 \\
Co &  -0.0632 &  -0.0248 & 1.382 & -0.028 & +0.014 \\
Cr &  +0.0708 &  +0.0422 & 1.420 & +0.033 & +0.101 \\
Cu &  -0.0122 &  -0.0067 & 1.413 & -0.002 & +0.027 \\
Dy &  -0.0265 &  -0.0584 & 2.023 & +0.051 & +0.001 \\
Er &  -0.0423 &  -0.0578 & 1.948 & +0.047 & -0.008 \\
Fe &  -0.0284 &  +0.0249 & 1.411 & +0.015 & +0.078 \\
Ga &   {---}  &   {---}  & 1.663 & -0.027 & +0.015 \\
Ge &   {---}  &  -0.1525 & 1.750 & -0.136 & +0.051 \\
Hf &  -0.0025 &  -0.0251 & 1.745 & -0.010 & +0.041 \\
Ir &  -0.0642 &  -0.0185 & 1.502 & -0.041 & +0.028 \\
La &  +0.0404 &  -0.0503 & 2.070 & +0.062 & +0.092 \\
Mg &  -0.0311 &  -0.0771 & 1.772 & -0.024 & -0.021 \\
Mn &  +0.0481 &  +0.0088 & 1.428 & +0.043 & +0.063 \\
Mo &  +0.0739 &   {---}  & 1.551 & +0.028 & +0.063 \\
Nb &  +0.0708 &  +0.0265 & 1.625 & +0.025 & +0.049 \\
Ni &  -0.0419 &  -0.0422 & 1.377 & -0.017 & +0.028 \\
Os &  -0.0471 &   {---}  & 1.493 & -0.051 & +0.019 \\
Pb &   {---}  &  -0.0326 & 1.932 & -0.009 & -0.031 \\
Pd &  -0.0271 &  -0.0204 & 1.517 & -0.005 & +0.033 \\
Pt &  -0.0275 &  -0.0074 & 1.534 & -0.003 & +0.031 \\
Re &  -0.0449 &   {---}  & 1.518 & -0.029 & +0.046 \\
Rh &  -0.0507 &  -0.0290 & 1.484 & -0.032 & +0.021 \\
Ru &  -0.0527 &  -0.0246 & 1.478 & -0.028 & +0.023 \\
Sc &  -0.0634 &  -0.0696 & 1.814 & -0.033 & -0.004 \\
Si &  -0.1806 &  -0.1338 & 1.683 & -0.211 & +0.048 \\
Sn &  -0.0530 &  -0.0563 & 1.863 & -0.058 & +0.004 \\
Ta &  +0.0576 &  -0.0073 & 1.626 & +0.022 & +0.037 \\
Tb &  -0.0256 &  -0.0510 & 1.973 & +0.044 & +0.004 \\
Ti &  -0.0162 &  -0.0370 & 1.615 & -0.005 & +0.032 \\
V  &  -0.0160 &  -0.0135 & 1.487 & +0.006 & +0.046 \\
W  &  +0.0615 &   {---}  & 1.558 & +0.016 & +0.064 \\
Y  &  -0.0435 &  -0.0614 & 1.990 & +0.007 & -0.006 \\
Zn &  -0.0297 &  -0.0434 & 1.536 & -0.020 & +0.007 \\
Zr &  +0.0003 &  -0.0276 & 1.771 & -0.001 & +0.033 \\
\end{longtable}

Key observations from Table~\ref{tab:delta}:
\begin{enumerate}
  \item Elements with large positive $\delta$ (Cr, Mo, Nb, Mn, W, Ta)
    cause positive $\Omega_\text{sf}$ deviations, i.e., lattice
    expansion beyond Vegard's law.
  \item Elements with large negative $\delta$ (Si, Ca, Co, Ir, Rh, Sc)
    cause lattice contraction.
  \item The $\delta^\text{B2}$ and $\delta^{\text{L1}_2}$ values are
    correlated but not identical, reflecting the different coordination
    environments.
  \item Minority-site radii show systematically larger deviations from
    pure-element values ($\langle|\Delta r_\text{min}|\rangle = 0.045$~\AA)
    than majority-site radii ($\langle|\Delta r_\text{maj}|\rangle =
    0.036$~\AA), consistent with the greater perturbation experienced
    in a fully heterogeneous coordination shell.
\end{enumerate}

\subsection{Additive Decomposition Quality}
\label{sec:additive_quality}

The additive model $\Omega_\text{sf} \approx \delta_A + \delta_B$
explains 62\% of the B2 variance ($R^2 = 0.62$, RMSE = 0.042) and
69\% of the L1$_2$ variance ($R^2 = 0.69$, RMSE = 0.029). The
remaining variance represents genuinely pairwise interactions (charge
transfer, $d$--$d$ hybridisation) that cannot be decomposed into
single-element contributions.

Despite this approximation, the additive model applied to HEA lattice
prediction achieves RMSE = 0.0215~\AA---only 0.0005~\AA{} (2.4\%)
worse than the full pairwise model (0.0210~\AA). This is because:
(a)~the pairwise residuals average out in multi-element sums; (b)~the
$\gamma$ optimisation partially absorbs systematic biases; (c)~the
additive model covers all 1,227 possible element pairs versus only
220 for pairwise data.

\begin{table}[htbp]
\centering
\caption{Comparison of pairwise versus additive $\Omega_\text{sf}$
  for HEA lattice prediction.}
\label{tab:pw_vs_add}
\begin{tabular}{@{}lcccc@{}}
\toprule
Approach & Parameters & Pairs covered & RMSE (\AA) & $\gamma$ \\
\midrule
Pairwise $\Omega_\text{sf}$ & 220 pairs & 220 & 0.0210 &
  $\gamma_\text{B} = 1.45$, $\gamma_\text{F} = 1.08$ \\
Additive $\delta_A + \delta_B$ & 39 elements & 1,227 & 0.0215 &
  $\gamma_\text{B} = -0.16$, $\gamma_\text{F} = -0.14$ \\
\bottomrule
\end{tabular}
\end{table}

\subsection{Structural Information in $\delta r$ vs.\ $\Omega_\text{sf}$}
\label{sec:structural}

\subsubsection{Mathematical proof: $\delta r$ is structure-invariant}

From Eq.~\eqref{eq:delta_r}, $\delta r = f(c_1, \ldots, c_n;\,
r_1, \ldots, r_n)$ is a function of composition $\{c_i\}$ and
pure-element radii $\{r_i\}$ only. No crystal structure argument
appears. For a binary system $A$--$B$ at any given composition
$c_A$, the value of $\delta r$ is:
\begin{equation}
  \delta r(c_A) = 100\,\frac{|r_A - r_B|\sqrt{c_A(1-c_A)}}
    {c_A r_A + (1-c_A) r_B}.
  \label{eq:delta_r_binary}
\end{equation}
This depends only on $c_A$, $r_A$, and $r_B$---not on whether the
alloy adopts L1$_2$ ($c_A = 0.25$ or $0.75$) or B2 ($c_A = 0.50$)
structure.

\subsubsection{Computational validation: 905 element pairs}

We computed $\delta r$ and $\Omega_\text{sf}$ for all 905 element
pairs having DFT data for L1$_2$ $A_3B$, $B_3A$, and B2 $AB$.

\paragraph{$\delta r$ shows zero structural scatter:}
Plotting $\delta r(A_3B)$ versus $\delta r(B_3A)$ yields a smooth,
one-dimensional curve with no scatter---confirming that $\delta r$ is
purely composition-determined. The same element pair at
$c = 0.25$ and $c = 0.75$ gives different $\delta r$ values, but
these differ only because the composition changed, not because the
structure changed.

\paragraph{$\Omega_\text{sf}$ shows large structural scatter:}
In contrast, $\Omega_\text{sf}(A_3B)$ versus $\Omega_\text{sf}(B_3A)$
shows a correlation of only $r = 0.15$ (near zero), demonstrating
that the same element pair produces very different volume deviations
depending on which element occupies the majority versus minority site.

\paragraph{L1$_2$ asymmetry:}
The mean absolute lattice constant difference between $A_3B$ and
$B_3A$ is 0.328~\AA, with a maximum of 2.74~\AA. This structural
asymmetry is invisible to $\delta r$ but fully captured by
$\Omega_\text{sf}$.

\subsection{Why Packing-Constrained Radii Are Insufficient}
\label{sec:packing_results}

We tested whether DFT-derived packing radii (Section~\ref{sec:packing})
can serve as an alternative. Fitting a single radius set $(r_A, r_B)$
to all three structures simultaneously yields:

\begin{table}[htbp]
\centering
\caption{Single packing radius RMSE for nearest-neighbour distance
  prediction.}
\label{tab:packing}
\begin{tabular}{@{}lc@{}}
\toprule
Structure & RMSE (d$_\text{nn}$, \AA) \\
\midrule
L1$_2$ $A_3B$ & 0.177 \\
L1$_2$ $B_3A$ & 0.202 \\
B2 $AB$        & 0.084 \\
Overall        & 0.162 \\
\bottomrule
\end{tabular}
\end{table}

The fundamental limitation is that $d_\text{nn} = r_A + r_B$ is a
\emph{symmetric} operation: it yields identical values for $A_3B$ and
$B_3A$, forcing $a(A_3B) = a(B_3A)$. In reality,
$\langle|a(A_3B) - a(B_3A)|\rangle = 0.328$~\AA. The packing model
cannot represent the role asymmetry (majority vs.\ minority) that is
essential for L1$_2$ compounds.

The volume-based approach (Section~\ref{sec:eff_radii}), by contrast,
allows $V_\text{maj}(A) \neq V_\text{min}(A)$, naturally encoding
the structural asymmetry without additional constraints.

\subsection{Multiphase Classification: $\delta r$ vs.\ $\delta_\text{sf}$}
\label{sec:phase}

We define a volume-size-factor-based mismatch:
\begin{equation}
  \delta_\text{sf} = 100\sqrt{
    \sum_i c_i \!\left(1 - \frac{1 + \bar{\Omega}_i}{1 +
    \bar{\Omega}}\right)^{\!2}},
  \label{eq:delta_sf}
\end{equation}
where $\bar{\Omega}_i = \sum_{j \ne i} c_j \Omega_\text{sf}(i,j)$.

On the 54-HEA multiphase database:

\begin{table}[htbp]
\centering
\caption{Phase classification performance (SS vs.\ non-SS).}
\label{tab:phase}
\begin{tabular}{@{}lccc@{}}
\toprule
Criterion & AUC & Accuracy & F1 \\
\midrule
$\delta r$ (threshold 4.72\%) & \textbf{0.873} & \textbf{0.800} & \textbf{0.766} \\
$\delta_\text{sf}$ (threshold 0.043) & 0.708 & 0.661 & 0.698 \\
Yang--Zhang ($\delta r < 6.6\%$, $\Omega > 1.1$) & --- & 0.661 & 0.732 \\
$\delta_\text{sf}$ + $\Omega$ combined & --- & 0.764 & 0.764 \\
\bottomrule
\end{tabular}
\end{table}

$\delta r$ outperforms $\delta_\text{sf}$ for phase classification
(AUC 0.873 vs.\ 0.708). This is an important negative result:
$\delta_\text{sf}$ encodes chemical bonding information, but
single-phase stability is primarily governed by geometric size
mismatch. The value of $\Omega_\text{sf}$ and $\delta_\text{sf}$
lies in \emph{lattice parameter prediction}, not phase classification.

\subsection{Machine Learning Residual Analysis}
\label{sec:ml_results}

Seven ML models were tested for residual correction. None improves
the RMSE by more than 0.001~\AA{} (Table~\ref{tab:comparison}).
Direct XGBoost on the 64 HEAs (23 features) yields RMSE = 0.123~\AA---
six times worse than Vegard's law---demonstrating the failure of
flexible ML on small datasets. Transfer learning from 1,062 binary
compounds produces RMSE = 0.032~\AA, worse than the physics baseline,
due to the large domain gap between ordered binaries and random HEA
solid solutions.


%% =====================================================================
\section{Discussion}
\label{sec:discussion}
%% =====================================================================

\subsection{Why $\delta r$ Cannot Capture Structural Information}

The inability of $\delta r$ to encode structural information is not an
empirical observation but a \emph{mathematical consequence} of its
definition. Since $\delta r = f(c_i, r_i)$ contains no structure
argument, it assigns identical values to $A_3B$ and $B_3A$ (at the
same composition) regardless of which element occupies the majority
site. The 905-pair analysis provides overwhelming statistical
confirmation: the $\delta r(A_3B)$--$\delta r(B_3A)$ plot shows zero
scatter, whereas the $\Omega_\text{sf}(A_3B)$--$\Omega_\text{sf}(B_3A)$
plot has correlation $r = 0.15$.

\subsection{Physical Origin of Structure-Dependent Volumes}

The structure dependence of effective atomic volumes arises from
electronic structure effects:
\begin{enumerate}
  \item \textbf{Coordination environment}: In L1$_2$ $A_3B$, the
    minority $B$ atom is surrounded by 12 $A$ neighbours (fully
    heterogeneous shell), while each $A$ atom has 8 same-type and
    4 opposite-type neighbours. This asymmetric coordination produces
    different charge distributions and atomic pressures.
  \item \textbf{Charge transfer}: Electronegativity differences drive
    charge redistribution that depends on the local composition
    (majority-site $A$ donates/accepts differently from
    minority-site $B$).
  \item \textbf{$d$--$d$ hybridisation}: Transition metal pairs exhibit
    directional bonding whose strength depends on the coordination
    geometry.
\end{enumerate}
These effects are absent in pure-element radii and in the packing
model (which treats atoms as hard spheres), but are naturally captured
by DFT atomic volumes.

\subsection{Limitations of Packing-Constrained Radii}

The packing model fails for a fundamental algebraic reason: the
nearest-neighbour distance $d_\text{nn} = r_A + r_B$ is symmetric
under exchange of majority and minority roles. Any model based on
$r_A + r_B$ will predict $a(A_3B) = a(B_3A)$, contradicting the
DFT data ($\langle|\Delta a|\rangle = 0.328$~\AA).

Structure-specific packing radii ($r_A^\text{maj}$,
$r_A^\text{min}$, $r_A^\text{B2}$) \emph{can} absorb structural
information, but this approach is functionally equivalent to the
volume-based method while introducing the additional, unphysical
assumption of rigid-sphere contact.

\subsection{Practical Implications of Additive Decomposition}

The additive model $\Omega_\text{sf} \approx \delta_A + \delta_B$
offers three practical advantages:
\begin{enumerate}
  \item \textbf{Compact representation}: 39 element parameters
    (Table~\ref{tab:delta}) versus 220 pairwise values.
  \item \textbf{Full coverage}: Any pair of the 39 elements can be
    predicted ($\binom{39}{2} = 741$ pairs), compared to only 220
    with measured pairwise data.
  \item \textbf{Reproducibility}: Another researcher can reproduce
    predictions using only Table~\ref{tab:delta}, without access to
    DFT databases.
\end{enumerate}

The 2.4\% RMSE penalty (0.0215 vs.\ 0.0210~\AA) is negligible for
practical screening, making the additive model the recommended
approach for rapid HEA lattice parameter estimation.

\subsection{Comparison with Alonso's Approach}

Our method differs from Alonso~\cite{Alonso2021} in three key aspects:
\begin{enumerate}
  \item \textbf{Data source}: DFT-computed lattice constants versus
    King's 1966 experimental measurements, eliminating reliance on
    decades-old data and extending to element pairs not covered by
    King.
  \item \textbf{Structure specificity}: B2 data for BCC HEAs and
    L1$_2$ data for FCC HEAs, versus structure-agnostic $\Omega_\text{sf}$.
  \item \textbf{Additive decomposition}: Element-level parameters
    enabling prediction for any element pair, versus pairwise-only
    coverage.
\end{enumerate}

\subsection{Noise Floor and Achievable Accuracy}

The noise floor $\sigma \approx 0.016$~\AA{} establishes a fundamental
limit. Our best RMSE (0.0210~\AA) is within a factor of 1.3 of this
limit, suggesting that further improvement requires either
higher-quality experimental data or models incorporating temperature
dependence and short-range order. The FCC RMSE of 0.0096~\AA{}
is \emph{below} the noise floor, indicating that the FCC model
predictions are limited by measurement uncertainty rather than model
accuracy.

\subsection{Negative Results and Their Significance}

Two significant negative results merit discussion:
\begin{enumerate}
  \item \textbf{ML residual correction failure}: Seven ML models
    fail to improve upon the physics model, confirming that the
    residuals are dominated by measurement noise rather than
    systematic model error.
  \item \textbf{$\delta_\text{sf}$ for phase classification}: The
    DFT-derived $\delta_\text{sf}$ performs \emph{worse} than simple
    $\delta r$ for single-phase prediction (AUC 0.708 vs.\ 0.873).
    This shows that chemical bonding effects encoded in
    $\Omega_\text{sf}$ are not directly relevant to
    solid-solution stability, which is governed primarily by
    geometric size compatibility.
\end{enumerate}


%% =====================================================================
\section{Conclusions}
\label{sec:conclusions}
%% =====================================================================

We have developed an improved framework for predicting HEA lattice
parameters that combines structure-specific DFT volume size factors
with element-wise additive decomposition. The principal findings are:

\begin{enumerate}
  \item Structure-specific DFT $\Omega_\text{sf}$ (B2 for BCC,
    L1$_2$ for FCC) with optimised $\gamma$ parameters achieves
    RMSE = 0.0210~\AA{} on 64 cubic HEAs, an 8.6\% improvement over
    Alonso's Eq.~10 (0.0229~\AA). FCC HEAs are predicted at
    0.0096~\AA, below the experimental noise floor.

  \item Element-wise additive decomposition
    ($\Omega_\text{sf} \approx \delta_A + \delta_B$) reduces 220
    pairwise parameters to 39 element-level values
    (Table~\ref{tab:delta}) with only 2.4\% RMSE penalty (0.0215
    vs.\ 0.0210~\AA), providing a practical, reproducible descriptor
    for rapid screening.

  \item Validation on 20 independent HEAs (RMSE = 0.0200~\AA)
    confirms generalisation capability.

  \item The conventional size mismatch $\delta r$ is mathematically
    incapable of encoding structural information, as proven by its
    definition ($\delta r = f(c_i, r_i)$, no structure argument)
    and validated on 905 binary pairs.

  \item Volume-derived effective radii naturally absorb structural
    effects through DFT electronic structure calculations, with
    minority-site perturbations 25\% larger than majority-site
    perturbations, consistent with the coordination environment
    asymmetry.

  \item Packing-constrained radii ($r_A + r_B = d_\text{nn}$) fail
    to capture L1$_2$ asymmetry due to the symmetric nature of the
    sum operation.

  \item ML residual corrections provide negligible improvement
    ($< 0.002$~\AA), confirming that the physics model captures
    the dominant variability and that residuals are noise-limited.
\end{enumerate}

Future work should focus on (a)~extending the VASP binary compound
database to increase pairwise coverage, (b)~incorporating
ternary/quaternary DFT data for multi-body corrections, and
(c)~temperature-dependent corrections for room-temperature lattice
parameters.

%% =====================================================================
%% References
%% =====================================================================

\begin{thebibliography}{99}

\bibitem{Yeh2004}
J.-W.~Yeh \emph{et al.}, Nanostructured high-entropy alloys with
multiple principal elements: novel alloy design concepts and outcomes,
\emph{Adv.\ Eng.\ Mater.}\ \textbf{6} (2004) 299--303.

\bibitem{Cantor2004}
B.~Cantor \emph{et al.}, Microstructural development in equiatomic
multicomponent alloys, \emph{Mater.\ Sci.\ Eng.\ A}\ \textbf{375--377}
(2004) 213--218.

\bibitem{Miracle2017}
D.B.~Miracle, O.N.~Senkov, A critical review of high entropy alloys and
related concepts, \emph{Acta Mater.}\ \textbf{122} (2017) 448--511.

\bibitem{Senkov2015}
O.N.~Senkov \emph{et al.}, Accelerated exploration of multi-principal
element alloys with solid solution phases,
\emph{Nat.\ Commun.}\ \textbf{6} (2015) 6529.

\bibitem{Vegard1921}
L.~Vegard, Die Konstitution der Mischkristalle und die Raumf\"ullung der
Atome, \emph{Z.\ Phys.}\ \textbf{5} (1921) 17--26.

\bibitem{King1966}
H.W.~King, Quantitative size-factors for metallic solid solutions,
\emph{J.\ Mater.\ Sci.}\ \textbf{1} (1966) 79--90.

\bibitem{Alonso2021}
J.A.~Alonso, Estimation of lattice parameters in cubic high-entropy
alloys, \emph{J.\ Phase Equilib.\ Diffus.}\ (2021).

\bibitem{Zhang2008}
Y.~Zhang \emph{et al.}, Solid-solution phase formation rules for
multi-component alloys,
\emph{Adv.\ Eng.\ Mater.}\ \textbf{10} (2008) 534--538.

\bibitem{Guo2011}
S.~Guo, C.T.~Liu, Phase stability in high entropy alloys: Formation of
solid-solution phase or amorphous phase,
\emph{Prog.\ Nat.\ Sci.: Mater.\ Int.}\ \textbf{21} (2011) 433--446.

\bibitem{Yang2012}
X.~Yang, Y.~Zhang, Prediction of high-entropy stabilized solid-solution
in multi-component alloys, \emph{Mater.\ Chem.\ Phys.}\ \textbf{132}
(2012) 233--238.

\bibitem{Jain2013}
A.~Jain \emph{et al.}, Commentary: The Materials Project: A materials
genome approach to accelerating materials innovation,
\emph{APL Mater.}\ \textbf{1} (2013) 011002.

\bibitem{Saal2013}
J.E.~Saal \emph{et al.}, Materials design and discovery with
high-throughput density functional theory: The Open Quantum Materials
Database (OQMD), \emph{JOM}\ \textbf{65} (2013) 1501--1509.

\bibitem{Chen2016}
T.~Chen, C.~Guestrin, XGBoost: A scalable tree boosting system,
\emph{Proc.\ 22nd ACM SIGKDD} (2016) 785--794.

\bibitem{Senkov2013}
O.N.~Senkov \emph{et al.}, Mechanical properties of
Nb$_{25}$Mo$_{25}$Ta$_{25}$W$_{25}$ and
V$_{20}$Nb$_{20}$Mo$_{20}$Ta$_{20}$W$_{20}$ refractory high entropy
alloys, \emph{Intermetallics}\ \textbf{19} (2011) 698--706.

\bibitem{Senkov2013mp}
O.N.~Senkov \emph{et al.}, Microstructure and room temperature properties
of a high-entropy TaNbHfZrTi alloy,
\emph{J.\ Alloys Compd.}\ \textbf{509} (2011) 6043--6048.

\bibitem{Yao2016}
H.W.~Yao \emph{et al.}, MoNbTaV medium-entropy alloy,
\emph{Entropy}\ \textbf{18} (2016) 189.

\bibitem{Otto2013}
F.~Otto \emph{et al.}, The influences of temperature and microstructure
on the tensile properties of a CoCrFeMnNi high-entropy alloy,
\emph{Acta Mater.}\ \textbf{61} (2013) 5743--5755.

\bibitem{Niu2017}
C.~Niu \emph{et al.}, Spin-driven ordering of Cr in the equiatomic high
entropy alloy NiFeCrCo, \emph{Sci.\ Rep.}\ \textbf{7} (2017) 7043.

\bibitem{Wang2019}
Z.~Wang \emph{et al.}, Atomic-size effect and solid solubility of
multicomponent alloys, \emph{Scr.\ Mater.}\ \textbf{94} (2015) 28--31.

\bibitem{Zaddach2013}
A.J.~Zaddach \emph{et al.}, Mechanical properties and stacking fault
energies of NiFeCrCoMn high-entropy alloy, \emph{JOM}\ \textbf{65}
(2013) 1780--1789.

\bibitem{Zhang2017nat}
Z.~Zhang \emph{et al.}, Dislocation mechanisms and 3D twin architectures
generate exceptional strength-ductility-toughness combination in CrCoNi
medium-entropy alloy, \emph{Nat.\ Commun.}\ \textbf{8} (2017) 14390.

\end{thebibliography}

\end{document}
